% =====================================================================
\section{Languages}
\label{sec:supp:languages}

Each of the six languages in \TT{Languages/} ships with a
\TT{Syntax} file (the AST), often an \TT{SOS} operational
semantics, and a \TT{Compile} file producing a
\DefRef{KernelGame}.  This section anchors at the compilation
\emph{shape} for each language; full Syntax content beyond the
top-level record is bunched.

\subsection{NFG: \texttt{Languages/NFG/Compile.lean}}
\label{sec:supp:nfg}\label{def:NFG.NFGGame}\label{def:NFG.toKernelGame}

\Anchor{NFG.NFGGame.toKernelGame}
\begin{leanbox}
structure NFGGame ($\iota$ : Type) (A : $\iota \to$ Type) where\\
\bd Outcome  : Type\\
\bd outcome  : ($\forall$ i, A i) $\to$ Outcome\\
\bd utility  : ($\forall$ i, A i) $\to$ $\iota \to \mathbb{R}$\\[3pt]
noncomputable def NFGGame.toKernelGame (G : NFGGame $\iota$ A) : KernelGame $\iota$ where\\
\bd Strategy := A\\
\bd Outcome  := G.Outcome\\
\bd utility  := fun $\omega$ i $\Rightarrow$ \dots\\
\bd outcomeKernel := fun $\sigma$ $\Rightarrow$ PMF.pure (G.outcome $\sigma$)
\end{leanbox}

\Mathematical
A normal-form game over a finite player set $\Players$ and
action sets $\{A_i\}_i$ is the textbook tuple
$(\Players, \{A_i\}_i, \util)$.  The compiler takes the strategy
type to be $A_i$, the outcome type to be a chosen $\Outcome$ with
a deterministic projection from joint actions, and the kernel to
be the Dirac evaluation $\Profile \mapsto \Dirac_{G.\mathrm{outcome}(\Profile)}$.
The compiled \KG\ is deterministic; randomness enters only at the
mixed extension.  The structure does not impose finiteness on the
action types $A_i$ — countably infinite action spaces are
admitted at the definitional level (see
\TT{Languages/NFG/CountableExample.lean} for an $\mathbb{N}$-valued
smoke test); the mixed-Nash existence theorem reintroduces
\TT{[Fintype (A i)]} where the proof requires it.

\Textbook
Osborne–Rubinstein Definition 11.1 (strategic-form game).  The
compilation is the identity on the strategic form modulo the
choice of outcome type.

\Faithfulness
Faithful, with reformulation:\ the textbook treats utility as a
direct function of profile; we route through an outcome type so
that a kernel-game-level theorem applies uniformly to NFGs and to
languages with non-trivial outcome types.

\Variants
\TT{Languages/NFG/} also packages worked subclasses, each as its
own \TT{NFGGame} construction:\ \TT{CongestionGame.lean}
(Rosenthal's potential), \TT{CoalitionalGame.lean} (the Shapley
value definition plus null-player, additivity, linearity, and
symmetry lemmas),
\TT{Bargaining.lean} (Nash bargaining), \TT{Matching.lean}
(stable-matching predicates), \TT{PublicGoods.lean},
\TT{Stackelberg.lean}.  Each subclass lives in its own namespace
and exposes its characteristic theorem.

\subsection{EFG: \texttt{Languages/EFG/}}
\label{sec:supp:efg}\label{def:EFG.InfoStructure}\label{def:EFG.GameTree}\label{def:EFG.EFGGame}\label{def:EFG.evalDist}

\Anchor{EFG.GameTree, EFG.evalDist, EFG.toKernelGame}
\begin{leanbox}
structure InfoStructure where\\
\bd n        : Nat\\
\bd Infoset  : Fin n $\to$ Type\\
\bd arity    : (p : Fin n) $\to$ Infoset p $\to$ Nat\\
\bd arity\_pos : $\forall$ p I, 0 < arity p I\\[3pt]
inductive GameTree (S : InfoStructure) (Outcome : Type) where\\
\bd $\mid$ terminal (z : Outcome)\\
\bd $\mid$ chance (k : Nat) ($\mu$ : PMF (Fin k)) (hk : 0 < k) (next : Fin k $\to$ GameTree S Outcome)\\
\bd $\mid$ decision \{p : S.Player\} (I : S.Infoset p) (next : S.Act I $\to$ GameTree S Outcome)\\[3pt]
structure EFGGame where\\
\bd inf : InfoStructure ; Outcome : Type\\
\bd tree : GameTree inf Outcome ; utility : Outcome $\to$ Payoff inf.Player\\[3pt]
noncomputable def GameTree.evalDist \{S : InfoStructure\} \{Outcome : Type\}\\
\bd ($\sigma$ : BehavioralProfile S) : GameTree S Outcome $\to$ PMF Outcome\\
\bd $\mid$ .terminal z $\Rightarrow$ PMF.pure z\\
\bd $\mid$ .chance \_ $\mu$ \_ next $\Rightarrow$ $\mu$.bind (fun b $\Rightarrow$ (next b).evalDist $\sigma$)\\
\bd $\mid$ .decision (p := p) I next $\Rightarrow$ ($\sigma$ p I).bind (fun a $\Rightarrow$ (next a).evalDist $\sigma$)
\end{leanbox}

\Mathematical
An extensive-form game is a triple $(\InfoStruct, \Tree, \util)$
where \InfoStruct\ records the players, their information sets,
and the action arity at each information set; the tree is
inductively built from terminal, chance, and decision nodes; and
the utility maps terminal outcomes to payoff vectors.  Strategies
come in three flavors:\ pure
$\Spure_p = \forall I, \mathrm{Fin}(\mathrm{arity}(p, I))$,
behavioral $\Sbeh_p = \forall I, \PMF{\mathrm{Fin}(\mathrm{arity}(p, I))}$,
and mixed $\Smix_p = \PMF{\Spure_p}$.  The function \TT{evalDist}
evaluates the tree under a behavioral profile to a $\PMF{\Outcome}$
by structural recursion:\ chance binds the chance distribution,
decision binds the player's behavioral distribution at the
relevant information set, terminal returns Dirac.  The compiled
kernel game has \TT{Strategy} the behavioral type, the EFG outcome
as outcome, \TT{evalDist} as kernel, and the EFG utility as utility.

\Textbook
Kuhn's representation~\citep{Kuhn1953}; Osborne–Rubinstein Chapter
6.  The \TT{InfoStructure} factoring through information-set IDs
(rather than directly through a player type) is a presentation
choice for the type theory:\ a decision node stores only its
information-set ID, and the player is recovered from the
information structure.  Augmented EFGs in the FOSG-paper
sense~\citep{Kovarik2022} are a thin wrapper adding public and
per-player history partitions (\TT{Languages/EFG/Augmented.lean}).

\Faithfulness
Faithful, with reformulation:\ distributional semantics rather
than deterministic-with-lift.  See chapter §3.2 for the
rationale.  Subgame-perfect equilibrium
(\TT{EFG.EFGGame.IsSubgamePerfectEq} in
\TT{Languages/EFG/Refinements.lean}) is defined as a Nash
equilibrium of every subgame; the ODP characterization
(\S\ref{sec:supp:odp}) makes this a one-shot deviation check.

\Variants
\TT{Languages/EFG/} contains \TT{Augmented.lean} (the augmented
EFG), \TT{SOS.lean} (operational semantics for the tree),
\TT{Compile.lean} and \TT{CompileObs.lean} (the latter compiling
to ObsModelCore for Kuhn), \TT{Kuhn.lean} (language-level Kuhn
specialization), \TT{Refinements.lean} (subgame perfection,
behavioral/pure refinements), \TT{Examples.lean}, \TT{Tests.lean}.

\subsection{MAID: \texttt{Languages/MAID/}}
\label{sec:supp:maid}\label{def:MAID.MAIDGame}

\Anchor{MAID.MAIDGame.toKernelGame}
\begin{leanbox}
structure MAIDGame where\\
\bd Players      : Type ; [fin : Fintype Players]\\
\bd Node         : Type\\
\bd parents      : Node $\to$ Finset Node\\
\bd nodeKind     : Node $\to$ NodeKind  -- chance, decision, utility\\
\bd domain       : Node $\to$ Type ; chanceDist, decisionInfo, utilityFn fields\\
\bd acyclic      : ...\\[3pt]
noncomputable def MAIDGame.toKernelGame (G : MAIDGame) : KernelGame G.Players
\end{leanbox}

\Mathematical
A multi-agent influence diagram~\citep{KollerMilch2003} is a
directed acyclic graph whose nodes are typed as chance, decision,
or utility, with edges encoding informational and causal
dependence.  The compiler implements a \emph{frontier evaluator}
that, at each step, picks a node whose parents are all assigned and
either samples its chance distribution, consults the player's
strategy at the relevant decision node, or accumulates its
utility.  Order-independence (the diamond lemma in
\TT{MAID/FrontierLemmas.lean}) makes the compilation
well-defined.

\Textbook
Koller–Milch Definition 1; Pfeffer–Gal Section 2.  The Lean
formalization restricts to finite-domain MAIDs (each node has a
\TT{Fintype} domain).  Strategic-relevance and d-separation
analyses, prominent in the MAID literature, are not formalized;
the order-independence content is.

\Faithfulness
Faithful, with restriction to finite-domain nodes.  See chapter
§5.4 for the rationale on why we accepted the proof obligations of
order-independent compilation rather than fixing a topological
order.

\Variants
\TT{Languages/MAID/} contains \TT{Syntax.lean}, \TT{SOS.lean},
\TT{FrontierEval.lean}, \TT{FrontierLemmas.lean},
\TT{FoldEval.lean}, \TT{Prefix.lean} (subdiagram framework),
\TT{Compile.lean}, \TT{CompileObs.lean}, \TT{Kuhn.lean},
\TT{Theorems.lean}.

\subsection{MultiRound: \texttt{Languages/MultiRound/}}
\label{sec:supp:mr}\label{def:MultiRoundGame}

\Anchor{MultiRoundGame, Round}
\begin{leanbox}
structure Round (n : Nat) (S V A Sig : Type) where\\
\bd signal     : S $\to$ PMF (Fin n $\to$ Sig)\\
\bd view       : Fin n $\to$ S $\to$ Sig $\to$ V\\
\bd transition : S $\to$ (Fin n $\to$ Option A) $\to$ S\\[3pt]
structure MultiRoundGame (n : Nat) (S V A Sig : Type) where\\
\bd init   : S\\
\bd rounds : List (Round n S V A Sig)\\[3pt]
noncomputable def MultiRoundGame.toKernelGame
\end{leanbox}

\Mathematical
A multi-round game is an initial state and a finite list of rounds
\citep{ShohamLeytonBrown2008}.  Each round samples a joint signal
from a state-dependent distribution, lets each player observe their
own state-and-signal-derived view, and applies a deterministic
transition keyed on the joint optional action.  Strategies are pure
and memoryless at the protocol level — a function from view to
optional action — with mixed/behavioral and recall structure
imposed externally.  The compilation generates the outcome kernel
by folding the round list and threading the state through the
chain of observations and transitions.

\Variants
\TT{Languages/MultiRound/RepeatedGame.lean} restricts to identical
rounds; \TT{StochasticGame.lean}~\citep{Shapley1953stochastic}
specializes the transition to be itself stochastic and
state-conditioned.  \TT{CompileObs.lean} and its linearizations
(\TT{CompileObsLin.lean}, \TT{CompileObsLinAdequacy.lean},
\TT{CompileObsLinRecall.lean}) compile to the
\DefRef{ObsModelCore} substrate that the Kuhn theorem
(\S\ref{sec:supp:kuhn}) is stated against.

\subsection{FOSG: \texttt{Languages/FOSG/}}
\label{sec:supp:fosg}\label{def:FOSG}

\Anchor{FOSG.toKernelGame}
\begin{leanbox}
structure FOSG ($\iota$ W A : Type) (Obs : $\iota \to$ Type) (Pub : Type) where\\
\bd init             : W\\
\bd active           : W $\to$ Finset $\iota$\\
\bd availableActions : (w : W) $\to \iota \to$ Set A\\
\bd terminal         : W $\to$ Prop\\
\bd transition       : ...\\
\bd reward, privObs, pubObs : ...\\
\bd terminal\_active\_eq\_empty : ...\\
\bd terminal\_no\_legal, nonterminal\_exists\_legal : ...\\[3pt]
noncomputable def FOSG.toKernelGame (G : FOSG \dots) [BoundedHorizon] : KernelGame $\iota$
\end{leanbox}

\Mathematical
A factored-observation stochastic game~\citep{Kovarik2022} is the
most general of the language layers:\ a state space, an active-set
function (the players who must move at each state), per-player
available actions, a transition kernel, a reward vector, a per-
player private observation, and a public observation.  Strategies
are functions of observation sequences; histories are recorded by
an explicit \TT{History} type.  Bounded horizon makes the
compilation finite.

\Textbook
Kovařík et al.~\citep{Kovarik2022} Definition 1.

\Variants
\TT{Languages/FOSG/} contains \TT{Basic.lean}, \TT{History.lean},
\TT{Information.lean}, \TT{Strategy.lean}, \TT{Values.lean},
\TT{Execution.lean}, \TT{OutcomeClosure.lean}, \TT{Serial.lean},
\TT{Compile.lean}, \TT{Kuhn.lean}, \TT{Theorems.lean}, with the
canonical \DefRef{ObsModelCore} bridge in \TT{Theorems/Kuhn/}.

\subsection{Intrinsic-form: \texttt{Languages/Intrinsic/}}
\label{sec:supp:intrinsic}

\Anchor{Intrinsic.productMixedToBehavioral}
\begin{leanbox}
structure WGame where\\
\bd Players, Agents : Type\\
\bd agentPlayer : Agents $\to$ Players\\
\bd Action : Agents $\to$ Type ; ...\\
\bd InfoStructure : ...\\[3pt]
noncomputable def productMixedToBehavioral (G : WGame) : ...\\
theorem behavioral\_realizes\_productMixed (G : WGame) : ...
\end{leanbox}

\Mathematical
The intrinsic (Witsenhausen-style) representation of a game
\citep{HeymannDeLaraChancelier2020} exposes agent decision
\emph{addresses} and information sets directly, without fixing an
order in which agents act.  Three strategy types are defined:\
\emph{mixed} (a distribution over a player's joint pure strategy
space $\prod_{a \in A_p} \Lambda_a$), \emph{product-mixed}
(independent randomization over each address), and
\emph{behavioral} (independent randomization at each information
set).  The unconditional product-mixed/behavioral equivalence is
the intrinsic-form analogue of one direction of Kuhn's theorem and
holds without any recall hypothesis.

\Textbook
Heymann–De Lara–Chancelier 2020 Definitions 8, 10, 11; Propositions
12, 13.  The intrinsic layer is structurally a separate syntactic
root than the EFG/MAID/MR/FOSG family.

% =====================================================================
\section{Bridges}
\label{sec:supp:bridges}

Bridges between languages are commuting-compilation theorems; the
expressiveness lens of \TT{Languages/Expressiveness/} packages
them in a uniform shape.

\subsection{Languages and reductions: anchor}
\label{sec:supp:expressiveness}\label{def:Expressiveness.Language}\label{def:Expressiveness.Reduction}

\Anchor{Expressiveness.Language, Expressiveness.Reduction}
\begin{leanbox}
namespace GameTheory.Languages.Expressiveness\\[1pt]
structure Language ($\iota$ : Type) where\\
\bd Syntax  : Type u\\
\bd compile : Syntax $\to$ KernelGame $\iota$\\[3pt]
structure Reduction \{$\iota$ : Type\} (L : Language $\iota$) (M : Language $\iota$)\\
\bd (R : KernelGame $\iota$ $\to$ KernelGame $\iota$ $\to$ Prop) where\\
\bd translate : L.Syntax $\to$ M.Syntax\\
\bd sound     : $\forall$ x : L.Syntax, R (L.compile x) (M.compile (translate x))
\end{leanbox}

\Mathematical
A \emph{language} for a fixed player type $\Players$ is a syntax
class together with a compilation $\compileTo : \mathrm{Syntax} \to
\KG_\Players$.  A \emph{reduction} from $L$ to $M$ over a semantic
relation $R$ is a translation on syntax together with a proof that
the compiled games are $R$-related for every source program.
Reductions compose under transitivity of $R$.

\Variants
The relation taxonomy in \TT{Expressiveness/Relations.lean}
catalogues:
$\mathrm{ProtocolEquivalent}$ (game-form-level isomorphism,
ignoring utility); $\mathrm{UtilityDistributionEquivalent}$
(bisimulation; both protocol and joint utility distribution
preserved); $\mathrm{UtilityDistributionSimulation}$ (the directed
version); $\mathrm{EUEquivalent}$ (equal expected utility for each
player at corresponding profiles); $\mathrm{EUSimulation}$ (directed
EU-comparison; equilibria of the source pull back to the target).
The bridges of the next subsection are reductions in this lens.

\subsection{Concrete bridges: bunched}
\label{sec:supp:concrete-bridges}

\paragraph{NFG → EFG (\TT{Languages/Bridges/NFG_EFG.lean}).}\
Each finite NFG is presented as a depth-1 EFG with a single
information set per player and a flat tree of $|\Act_i|$ branches
at each.  Soundness at the EU-isomorphism level.

\paragraph{NFG → FOSG (\TT{Languages/Bridges/NFG_FOSG.lean}).}\
Each finite NFG is presented as a one-shot FOSG with two world
states ($\mathsf{init}, \mathsf{terminal}$), every player active
at $\mathsf{init}$, no observations, and a single deterministic
transition delivering the NFG payoff as the FOSG reward.
Soundness:\
\TT{NFGGame.toFOSG_terminal_iff} and the corresponding EU
equivalence.

\paragraph{EFG → NFG (\TT{Languages/Bridges/EFG_NFG.lean}).}\
The strategic form of a finite EFG:\ pure strategy = function from
information set to action, payoff = expected payoff under the
behavioral lift of that pure strategy.  Soundness at the
EU-equivalence level.

\paragraph{MAID → EFG (\TT{Languages/Bridges/MAID_EFG.lean}).}\
A bounded MAID compiles to an EFG via a topological scheduling of
its decision nodes; soundness at the strategic-form refinement
level (information sets preserved).

\paragraph{EFG → FOSG (\TT{Languages/Bridges/FOSG/}).}\
\TT{ObsModelCore.lean}, \TT{AugmentedEFG.lean},
\TT{SerialExec.lean}.  An EFG with serial execution is presented
as an FOSG.  Soundness at the game-form simulation level.

\subsection{Solution-concept transport}
\label{sec:supp:transport}

The main chapter's map-based transport theorem
(\TT{Core/UtilityInvariance.lean}) states that a protocol map with
bijective \TT{stratMap} and bijective \TT{outcomeMap} transports the
preference-parameterized Nash, best-response, dominance,
weak/strict dominance, correlated-equilibrium, and
coarse-correlated-equilibrium predicates pointwise across the
bijection on profiles.  Specializing to a kernel-game simulation
(utility-preserving) yields the EU-specific transport laws.

The observed-law transport layer
(\TT{Concepts/DeviationSimulation.lean}) removes the need for an
outcome bijection.  It compares two games through views into a
common observed carrier $\Omega$, relates profiles or profile
distributions by equality of the pushed-forward outcome law, and
requires one-way simulation of target deviations by source
deviations.  Under these hypotheses, Nash and arbitrary
\TT{ProfileDeviationFamily} equilibria transfer from source to
target; a two-way simulation gives an iff.  These results are the
operational content of bridges whose target presentation contains
extra trace data or serialized administrative structure:\ a Nash
equilibrium of the abstract source remains Nash after compilation
provided all compiled unilateral deviations are observed already by
some source-side unilateral deviation.

% =====================================================================
\section{Major theorems}
\label{sec:supp:theorems}

\subsection{Mixed Nash existence: \texttt{KernelGame.mixed\_nash\_exists\_of\_bounded}}
\label{sec:supp:nash-exists}\label{def:mixed_nash_exists}

\Anchor{mixed_nash_exists_of_bounded, mixed_nash_exists}
\begin{leanbox}
theorem KernelGame.mixed\_nash\_exists\_of\_bounded (G : KernelGame $\iota$)\\
\bd [Fintype $\iota$] [DecidableEq $\iota$]\\
\bd [$\forall$ i, Finite (G.Strategy i)] [$\forall$ i, Nonempty (G.Strategy i)]\\
\bd \{$C$ : $\iota \to \mathbb{R}$\} (hbd : $\forall$ who $\omega$, |G.utility $\omega$ who| $\le$ C who) :\\
\bd $\exists$ $\sigma$ : $\forall$ i, PMF (G.Strategy i),\ G.mixedExtension.IsNash $\sigma$\\[3pt]
theorem KernelGame.mixed\_nash\_exists (G : KernelGame $\iota$)\\
\bd [Fintype $\iota$] [DecidableEq $\iota$]\\
\bd [$\forall$ i, Finite (G.Strategy i)] [$\forall$ i, Nonempty (G.Strategy i)]\\
\bd [Finite G.Outcome] :\\
\bd $\exists$ $\sigma$ : $\forall$ i, PMF (G.Strategy i),\ G.mixedExtension.IsNash $\sigma$
\end{leanbox}

\Mathematical
Every finite-strategy game with bounded player utilities has a
mixed Nash equilibrium.  Hypotheses:\ finite player type,
decidable equality on players, finite non-empty strategy space per
player, and a playerwise bound $|\util(\omega,i)| \le C_i$.
The outcome carrier need not be finite.  The theorem
\TT{mixed_nash_exists} is the finite-outcome wrapper that derives
boundedness automatically.

\paragraph{Proof structure.}\
The proof follows Nash's gain-map argument~\citep{Nash1951}.
Define
$\mathrm{gain}_i(\Profile, s) = \EU(\KG^{\mathrm{mix}}, \Profile[i \mapsto \Dirac_s], i) - \EU(\KG^{\mathrm{mix}}, \Profile, i)$,
the gain to player $i$ from deviating to pure $s$, and
\[
  \Phi_i(\Profile)(s) \;=\;
    \frac{\Profile_i(s) + (\mathrm{gain}_i(\Profile, s))^+}
         {1 + \sum_{s'} (\mathrm{gain}_i(\Profile, s'))^+}.
\]
$\Phi$ is a continuous self-map of the product simplex
$\prod_i \Delta(\Strat_i)$.  Brouwer's fixed-point theorem (mathlib's
\TT{brouwerFixedPoint} on a closed disk, lifted to the product
simplex via the homeomorphism in
\TT{Concepts/ProductSimplexBrouwer.lean}) gives a fixed point
$\Profile^*$.  The algebraic step in
\TT{Theorems/NashExistenceMixed.lean} (section
\TT{NashMapAlgebra}) shows that a fixed point of $\Phi$ is a Nash
equilibrium:\ the positive-part bracketing forces every player's
gain to vanish at $\Profile^*$.

\Textbook
Nash 1951 Theorem 1; Osborne–Rubinstein Proposition 33.1.

\Faithfulness
Faithful.  The library is honest about non-constructivity:\
\TT{mixed_nash_exists_of_bounded} is \TT{noncomputable} because Brouwer is.
A constructive route via Sperner's lemma is listed in chapter
§10 as a future direction.

\subsection{Correlated and coarse correlated equilibrium existence}
\label{sec:supp:ce-exists}\label{def:correlatedEq_exists}

\begin{leanbox}
theorem KernelGame.correlatedEq\_exists\_of\_bounded       (G : KernelGame $\iota$) [...]
theorem KernelGame.correlatedEq\_exists                   (G : KernelGame $\iota$) [...]
theorem KernelGame.coarseCorrelatedEq\_exists\_of\_bounded (G : KernelGame $\iota$) [...]
theorem KernelGame.coarseCorrelatedEq\_exists             (G : KernelGame $\iota$) [...]
\end{leanbox}

The independent product
$\bigotimes_i \Profile^*_i$ over the mixed Nash profile supplied by
the main chapter's mixed-existence theorem is a correlated equilibrium
of the original \KG;\ this is the route taken by
\TT{Theorems/CorrelatedEqExistence.lean}.  CCE existence is the
corollary $\Pref$-CE $\subseteq$ $\Pref$-CCE.  See chapter §6.2
for the alternative no-regret-learning route, not formalized.  The
\TT{\_of\_bounded} statements require finite strategy spaces but
not finite outcomes.

\subsection{Minimax theorem: \texttt{KernelGame.von\_neumann\_minimax\_of\_bounded}}
\label{sec:supp:minimax}\label{def:von_neumann_minimax}

\Anchor{von_neumann_minimax_of_bounded, von_neumann_minimax}
\begin{leanbox}
theorem KernelGame.von\_neumann\_minimax\_of\_bounded (G : KernelGame (Fin 2))\\
\bd [$\forall$ i, Finite (G.Strategy i)] [$\forall$ i, Nonempty (G.Strategy i)]\\
\bd (hzs : G.IsZeroSum)\\
\bd \{$C$ : Fin 2 $\to \mathbb{R}$\} (hbd : $\forall$ i $\omega$, |G.utility $\omega$ i| $\le$ C i) :\\
\bd $\exists$ (v : $\mathbb{R}$) ($\sigma$ : Profile G.mixedExtension),\\
\bdd G.mixedExtension.IsNash $\sigma$ $\wedge$ G.mixedExtension.eu $\sigma$ 0 = v $\wedge$\\
\bdd ($\forall$ s$_1$, G.mixedExtension.eu (Function.update $\sigma$ 1 s$_1$) 0 $\ge$ v) $\wedge$\\
\bdd ($\forall$ s$_0$, G.mixedExtension.eu (Function.update $\sigma$ 0 s$_0$) 0 $\le$ v)\\[3pt]
theorem KernelGame.von\_neumann\_minimax (G : KernelGame (Fin 2))\\
\bd [Finite G.Outcome] [$\forall$ i, Finite (G.Strategy i)] [$\forall$ i, Nonempty (G.Strategy i)]\\
\bd (hzs : G.IsZeroSum) :\\
\bd $\exists$ (v : $\mathbb{R}$) ($\sigma$ : Profile G.mixedExtension),\\
\bdd G.mixedExtension.IsNash $\sigma$ $\wedge$ G.mixedExtension.eu $\sigma$ 0 = v $\wedge$\\
\bdd ($\forall$ s$_1$, G.mixedExtension.eu (Function.update $\sigma$ 1 s$_1$) 0 $\ge$ v) $\wedge$\\
\bdd ($\forall$ s$_0$, G.mixedExtension.eu (Function.update $\sigma$ 0 s$_0$) 0 $\le$ v)
\end{leanbox}

\Mathematical
Every finite-strategy two-player zero-sum game with bounded utility
has a mixed Nash profile whose player-0 payoff is a saddle value:\
player 1 cannot lower it by deviating and player 0 cannot raise it
by deviating.  A separate lemma records that all mixed Nash equilibria yield the same
expected utility for player 0.  The usual finite-game
max-min/min-max equality is the textbook reading of these saddle
inequalities, but it is not the literal statement packaged here.

\paragraph{Proof structure.}\
Existence of a mixed Nash profile is \TT{mixed_nash_exists_of_bounded}.
Equality of EU at any two Nash profiles is the saddle-point
property of zero-sum Nash, recorded as
\TT{IsZeroSum.nash_eu_eq_of_bounded} in \TT{Theorems/Minimax.lean}.

\Textbook
von Neumann 1928~\citep{vN1928}; Maschler–Solan–Zamir Theorem 7.6.

\subsection{Backward induction (Zermelo) and the ODP}
\label{sec:supp:odp}\label{def:zermelo}\label{def:oneShotDeviation_iff_spe}\label{def:HasNoOneShotDeviation}

\Anchor{HasNoOneShotDeviation, oneShotDeviation_iff_spe, zermelo}
\begin{leanbox}
def HasNoOneShotDeviation (G : EFGGame)\\
\bd ($\sigma$ : PureProfile G.inf) : Prop :=\\
\bd $\forall$ \{p\} (I : G.inf.Infoset p) (next : G.inf.Act I $\to$ GameTree G.inf G.Outcome),\\
\bdd ($\exists$ h, ReachBy h G.tree (.decision I next)) $\to$\\
\bdd $\forall$ a' : G.inf.Act I,\\
\bdd \bd expect ((next ($\sigma$ p I)).evalDist (pureToBehavioral $\sigma$)) (G.utility · p) $\ge$\\
\bdd \bd expect ((next a').evalDist (pureToBehavioral $\sigma$)) (G.utility · p)\\[3pt]
theorem oneShotDeviation\_iff\_spe\_of\_bounded (G : EFGGame)\\
\bd ($\sigma$ : PureProfile G.inf) (hpi : IsPerfectInfo G.tree)\\
\bd \{$C$ : G.inf.Player $\to \mathbb{R}$\} (hbd : $\forall$ p $\omega$, |G.utility $\omega$ p| $\le$ C p) :\\
\bd G.IsSubgamePerfectEq $\sigma$ $\leftrightarrow$ HasNoOneShotDeviation G $\sigma$\\[3pt]
theorem zermelo\_of\_bounded (G : EFGGame) (hpi : IsPerfectInfo G.tree)\\
\bd \{$C$ : G.inf.Player $\to \mathbb{R}$\} (hbd : $\forall$ p $\omega$, |G.utility $\omega$ p| $\le$ C p) :\\
\bd $\exists$ $\sigma$ : PureProfile G.inf, G.IsSubgamePerfectEq $\sigma$
\end{leanbox}

\Mathematical
The one-shot deviation principle:\ in a finite perfect-information
EFG, a pure profile is subgame-perfect iff it has no profitable
one-shot deviation at any reachable decision node.  Zermelo's
theorem:\ every finite perfect-information EFG has a pure
subgame-perfect equilibrium, constructed by backward induction.
The displayed theorems are the bounded-utility forms; finite
outcome wrappers \TT{oneShotDeviation\_iff\_spe} and \TT{zermelo}
derive boundedness automatically.

\paragraph{Proof structure.}\
\TT{Theorems/Zermelo.lean}.
\TT{perfectInfo_disjoint_subtrees}:\ in a perfect-info tree,
sibling subtrees have disjoint information-set supports, so
backward induction can fix actions at one decision node without
affecting optimality elsewhere.  \TT{exists_noOSD}:\ the inductive
backward construction yields a profile with no profitable one-shot
deviation.  Zermelo follows from \TT{exists_noOSD} and the ODP.

\Textbook
Zermelo 1913 (the chess theorem); modern treatment in
Selten 1965~\citep{Selten1965, Selten1975}; Osborne–Rubinstein
Chapter 6.4.  ODP is folklore (see Fudenberg–Tirole §3.3 for a
modern statement).

\subsection{Kuhn's theorem: \texttt{Theorems/Kuhn/}}
\label{sec:supp:kuhn}\label{def:ObsModelCore}\label{def:kuhn_behavioral_to_mixed}\label{def:kuhn_mixed_to_behavioral_semantic}

\Anchor{ObsModelCore, kuhn_behavioral_to_mixed, kuhn_mixed_to_behavioral_semantic}
\begin{leanbox}
structure ObsModelCore ($\iota\,\sigma$ : Type) (Obs : $\iota \to$ Type)\\
\bd (Act : (i : $\iota$) $\to$ Obs i $\to$ Type) where\\
\bd  -- world state, transition kernel, per-player observation, ...\\[3pt]
def NoNontrivialInfoStateRepeat (O : ObsModelCore ...) : Prop\\
def HorizonSeparation                (O : ObsModelCore ...) : Prop\\
def ActionPosteriorLocal            (O : ObsModelCore ...) : Prop\\
def PerStepPlayerRecall             (O : ObsModelCore ...) : Prop\\[3pt]
theorem ObsModelCore.kuhn\_behavioral\_to\_mixed\\
\bd (O : ObsModelCore ...) [...] (hH : HorizonSeparation O) (hR : NoNontrivialInfoStateRepeat O) :\\
\bd $\forall$ $\beta$ : BehavioralProfile O,\ runDist O $\beta$ = (behavioralToMixedJoint $\beta$).bind (runDistPure O)\\[3pt]
theorem ObsModelCore.kuhn\_mixed\_to\_behavioral\_semantic\\
\bd (O : ObsModelCore ...) [...] (hL : ActionPosteriorLocal O) :\\
\bd $\forall$ $\mu$ : PMF (PureProfile O), $\exists$ $\beta$, runDist O $\beta$ = $\mu$.bind (runDistPure O)
\end{leanbox}

\Mathematical
The library proves Kuhn's behavioral$\,\leftrightarrow\,$mixed
equivalence at a representation-neutral substrate, the
\TT{ObsModelCore} record:\ a tuple recording the world state, each
player's observation type, the per-observation action arity, and
a transition kernel.  Behavioral profiles are per-player
$\InfoState_i \to \PMF{\Act_i}$; mixed profiles are per-player
distributions over the deterministic functions
$\InfoState_i \to \Act_i$.

The four hypothesis predicates control which direction of Kuhn
applies:
\begin{itemize}\itemsep1pt
\item \emph{NoNontrivialInfoStateRepeat}:\ no information state
  is visited twice along any trace (except at trivially-arity
  states).  This enforces tree-shaped reachability and is needed
  for the $B \to M$ direction's product-of-information-states
  construction to be well-defined.
\item \emph{HorizonSeparation}:\ a sequential separation of
  histories by depth, also needed for $B \to M$.
\item \emph{ActionPosteriorLocal}:\ the conditional distribution
  of an action given an information state does not depend on
  information outside the player's recall.  Sufficient for $M \to
  B$.  Strictly weaker than perfect recall in the textbook sense.
\item \emph{PerStepPlayerRecall}:\ the syntactic version of
  ActionPosteriorLocal that is automatic for snapshot-refined
  information states.
\end{itemize}

The $B \to M$ direction (\TT{kuhn_behavioral_to_mixed}) constructs
the equivalent mixed profile by the independent product of the
behavioral distributions across information states; the
$M \to B$ direction (\TT{kuhn_mixed_to_behavioral_semantic})
constructs the behavioral profile by extracting the
information-state-conditional distributions from $\mu$.

\Textbook
Kuhn 1953 Theorem 4 (perfect-recall version); Aumann 1964 (recall
hypothesis sharpening); Selten 1975 §6 (textbook treatment).

\Faithfulness
Strict generalization.  The textbook Kuhn's theorem requires
perfect recall; the library proves both directions under
strictly weaker hypotheses.  The $B \to M$ direction needs no
recall at all; the $M \to B$ direction needs only
ActionPosteriorLocal, which holds whenever the information
structure has the per-step locality property — a condition
strictly weaker than full perfect recall.  See chapter §6.5 for
the rationale.

\Variants
\TT{Theorems/Kuhn/} contains the proof split:\
\TT{ObsModel.lean} (the snapshot-refined wrapper),
\TT{KuhnModel.lean} (alternative carrier),
\TT{BehavioralToMixedCore.lean}
(the $B \to M$ proof, ~613 lines),
\TT{MixedToBehavioralCore.lean} (the $M \to B$ proof, ~951 lines),
\TT{CorrelatedRealization.lean} (an alternative
$M \to B$ proof via the correlated-equilibrium realization formula).
The package also exposes generic outcome-equality schemas
(\TT{KuhnBehavioralToMixedOutcome},
\TT{KuhnMixedToBehavioralViaOutcome},
\TT{KuhnCompleteViaOutcome} in \TT{Theorems/Kuhn.lean}) for
language-specific Kuhn instantiations.

% =====================================================================
\section{Mechanism design and auctions}
\label{sec:supp:mechanism}

\subsection{Bayesian games: \texttt{Mechanism/BayesianGame.lean}}
\label{sec:supp:bayesian}\label{def:BayesianGame}\label{def:BayesianGame.toKernelGame}

\Anchor{BayesianGame, BayesianGame.toKernelGame}
\begin{leanbox}
structure BayesianGame ($\iota$ : Type) [Fintype $\iota$] where\\
\bd $\Theta$       : $\iota \to$ Type ; [Fintype, Nonempty]\\
\bd Act     : $\iota \to$ Type ; [Fintype, Nonempty]\\
\bd prior   : PMF ($\forall$ i, $\Theta$ i)\\
\bd utility : ($\forall$ i, $\Theta$ i) $\to$ ($\forall$ i, Act i) $\to \iota \to \mathbb{R}$\\[3pt]
abbrev Strategy (B : BayesianGame $\iota$) (i : $\iota$) := B.$\Theta$ i $\to$ B.Act i\\[3pt]
noncomputable def exAnteEU (B : BayesianGame $\iota$) ($\sigma$ : B.BProfile) (who : $\iota$) : $\mathbb{R}$ :=\\
\bd expect B.prior (fun $\theta$ $\Rightarrow$ B.utility $\theta$ (fun i $\Rightarrow$ $\sigma$ i ($\theta$ i)) who)\\[3pt]
noncomputable def toKernelGame (B : BayesianGame $\iota$) : KernelGame $\iota$ :=\\
\bd KernelGame.ofEU B.Strategy (fun $\sigma$ i $\Rightarrow$ B.exAnteEU $\sigma$ i)\\[3pt]
def BayesNash (B : BayesianGame $\iota$) ($\sigma$ : B.BProfile) : Prop :=\\
\bd B.toKernelGame.IsNash $\sigma$\\[3pt]
theorem bayesNash\_iff\_exAnteEU (B : BayesianGame $\iota$) ($\sigma$ : B.BProfile) :\\
\bd B.BayesNash $\sigma$ $\leftrightarrow$ $\forall$ (who) (s'), B.exAnteEU $\sigma$ who $\ge$ B.exAnteEU (Function.update $\sigma$ who s') who
\end{leanbox}

\Mathematical
A Bayesian game is a tuple
$(\Theta, \mathrm{Act}, \mu, u)$ with $\Theta_i$ the type space of
player $i$, $\mathrm{Act}_i$ the action space, $\mu \in
\PMF{\prod_i \Theta_i}$ the common prior, and $u$ the
type-and-action-dependent payoff.  A strategy is a function
$\Theta_i \to \mathrm{Act}_i$.  The compilation goes through
\TT{KernelGame.ofEU}:\ the strategy type is preserved, the outcome
type is $\Players \to \RR$, and the kernel is the Dirac evaluation
at the ex-ante expected utility \TT{exAnteEU}.  Bayes-Nash
equilibrium is exactly Nash on the compiled kernel game,
characterized concretely by \TT{bayesNash_iff_exAnteEU}.

\Textbook
Harsanyi 1967~\citep{Harsanyi1967} type-agent representation;
Osborne–Rubinstein Chapter 24.

\Faithfulness
Reformulation.  The textbook Bayesian Nash equilibrium quantifies
over type-conditional deviations; we use the equivalent ex-ante
formulation, where a deviation is a function $\Theta_i \to
\mathrm{Act}_i$ replacing the player's strategy globally.  The
ex-post and interim variants are derived definitions.
\emph{Restriction:}\ finite type and action spaces.

\subsection{General mechanisms and the revelation principle}
\label{sec:supp:revelation}\label{def:GeneralMechanism}\label{def:revelation_principle}

\Anchor{GeneralMechanism, revelation_principle}
\begin{leanbox}
structure GeneralMechanism ($\iota$ : Type) [Fintype $\iota$] where\\
\bd $\Theta$       : $\iota \to$ Type ; Act : $\iota \to$ Type\\
\bd outcome : ($\forall$ i, Act i) $\to \iota \to \mathbb{R}$\\[3pt]
def StrategyProfile (M : GeneralMechanism $\iota$) := $\forall$ i, M.$\Theta$ i $\to$ M.Act i\\[3pt]
def IsBNE (M : GeneralMechanism $\iota$) ($\mu$ : PMF ($\forall$ i, M.$\Theta$ i)) ($\sigma$ : M.StrategyProfile) : Prop\\[3pt]
noncomputable def toDirect (M : GeneralMechanism $\iota$) ($\sigma$ : M.StrategyProfile) : Mechanism $\iota$\\[3pt]
theorem revelation\_principle (M : GeneralMechanism $\iota$) ($\mu$ : PMF ($\forall$ i, M.$\Theta$ i))\\
\bd ($\sigma$ : M.StrategyProfile) (h : M.IsBNE $\mu$ $\sigma$) :\\
\bd (M.toDirect $\sigma$).isBIC $\mu$
\end{leanbox}

\Mathematical
A general mechanism has separate type spaces and action spaces;
strategies map types to actions.  A direct mechanism has
$\mathrm{Act}_i = \Theta_i$.  Bayesian incentive compatibility (BIC)
of a direct mechanism is BNE under truthful reporting.  The
revelation principle says:\ for any general mechanism with a BNE
$\sigma$, the direct mechanism that composes reports with $\sigma$
is BIC in the library's constant-misreport sense.

\Textbook
Myerson 1981~\citep{Myerson1981}; Maschler–Solan–Zamir Chapter
17.

\Faithfulness
Faithful to the finite direct-mechanism construction, with one
restriction:\ \TT{Mechanism.isBIC} quantifies over constant
misreports $\theta'_i$, whereas the stronger Bayes-Nash predicate
over the induced Bayesian game quantifies over arbitrary
type-dependent reporting strategies.  The library includes
\TT{isBIC_of_truthful_bayesNash} to relate the two levels.

\subsection{Auctions: bunched}
\label{sec:supp:auctions}\label{def:vickrey_truthful_dominant}

\Anchor{vickrey_truthful_dominant}
\begin{leanbox}
noncomputable def vickreyPayoff (v : Fin n $\to \mathbb{R}$) (bids : Fin n $\to \mathbb{R}$) (who : Fin n) : $\mathbb{R}$ :=\\
\bd if bids who > maxOtherBid bids who then v who - maxOtherBid bids who else 0\\[3pt]
theorem vickrey\_truthful\_dominant (v : Fin n $\to \mathbb{R}$) (who : Fin n) :\\
\bd $\forall$ bids b', vickreyPayoff v (Function.update bids who (v who)) who $\ge$\\
\bdd vickreyPayoff v (Function.update bids who b') who
\end{leanbox}

\Mathematical
In a second-price (Vickrey) auction with valuations $v$, the
payoff is $v_i - \max_{j \neq i} b_j$ for the highest bidder
($b_i$) if their bid strictly exceeds the rest, and $0$ otherwise.
Truthful bidding $b_i = v_i$ is a weakly dominant strategy:\ for
any opponents' bids and any alternative bid $b'$, the payoff under
truthful bidding is at least as large.

\Textbook
Vickrey 1961~\citep{Vickrey1961}.

\Faithfulness
Faithful.  The proof is a case analysis on whether $b'$ is above
or below $\max_{j \neq i} b_j$, in each case verifying that the
payoff under truthful bidding matches or exceeds.

\Variants
\TT{Auctions/FirstPrice.lean}, \TT{Auctions/AllPay.lean},
\TT{Auctions/VCG.lean}.  First-price constructs the deterministic
first-price payoff game and proves that no single bid is dominant.
All-pay records payoff facts such as overbidding being unprofitable
and rent dissipation.  VCG packages the efficient-allocation setup
and proves truthful reporting is dominant
(Clarke 1971, Groves 1973).

\subsection{Social choice}
\label{sec:supp:social-choice}

\TT{Mechanism/SocialChoice.lean} formalizes a finite set of
alternatives, a finite set of voters with strict preferences over
alternatives, and structural notions:\ preference profiles,
utilitarian aggregation, the existence of a welfare-maximizing
alternative.  The classical impossibility theorems (Arrow,
Gibbard–Satterthwaite, Sen) are not formalized; see chapter §10.

% =====================================================================
\section{Worked examples}
\label{sec:supp:examples}

The chapter's §6 (Worked example) traces Prisoner's Dilemma at a
high level.  This section gives the full Lean text and the
unfolding of the relevant predicates.

\subsection{Prisoner's Dilemma}
\label{sec:supp:examples:pd}

\begin{leanbox}
namespace NFG\\[1pt]
inductive PDAction where $\mid$ cooperate $\mid$ defect\\
deriving DecidableEq, Repr\\[3pt]
def prisonersDilemma : NFGGame Bool (fun \_ $\Rightarrow$ PDAction) where\\
\bd Outcome := $\forall$ \_ : Bool, PDAction\\
\bd outcome := id\\
\bd utility s p :=\\
\bdd match s true, s false, p with\\
\bdd $\mid$ defect, defect, true $\Rightarrow$ 1\\
\bdd $\mid$ cooperate, defect, true $\Rightarrow$ 0\\
\bdd $\mid$ defect, cooperate, true $\Rightarrow$ 3\\
\bdd $\mid$ cooperate, cooperate, true $\Rightarrow$ 2\\
\bdd $\mid$ ...\\[3pt]
def pd\_defect\_defect : StrategyProfile (fun \_ : Bool $\Rightarrow$ PDAction) :=\\
\bd fun \_ $\Rightarrow$ PDAction.defect\\[3pt]
theorem pd\_defect\_is\_nash : IsNashPure prisonersDilemma pd\_defect\_defect := by\\
\bd intro i a' ; cases i $<;>$ cases a' $<;>$\\
\bdd simp [prisonersDilemma, pd\_defect\_defect, deviate, Function.update]
\end{leanbox}

\paragraph{Faithfulness check.}\
The predicate \TT{IsNashPure prisonersDilemma pd_defect_defect}
unfolds (after \TT{intro i a'}) to four cases, one per
$(i \in \{T, F\}, a' \in \{C, D\})$ pair:
\begin{itemize}\itemsep1pt
\item $i = T, a' = C$:\ payoff at $(D, D)$ for player $T$ is $1$,
  payoff at $(C, D)$ for player $T$ is $0$.  $1 \ge 0$. \checkmark
\item $i = T, a' = D$:\ same profile.  $1 \ge 1$. \checkmark
\item $i = F, a' = C$:\ payoff at $(D, D)$ for $F$ is $1$, payoff
  at $(D, C)$ for $F$ is $0$.  $1 \ge 0$. \checkmark
\item $i = F, a' = D$:\ same profile.  $1 \ge 1$. \checkmark
\end{itemize}
Each inequality is the textbook Nash check.  The Lean proof closes
all four with one \TT{simp}.

The companion theorem \TT{pd_coop_not_nash} :\
\TT{¬ IsNashPure prisonersDilemma pd_coop_coop}
exhibits a profitable deviation:\ player $T$ deviates from $C$ to
$D$ and gains payoff $3 - 2 = 1$.  The Lean proof extracts the
$T$, $D$ instance of $\IsNash$, simplifies, and closes with
\TT{linarith}.

\subsection{Matching Pennies}
\label{sec:supp:examples:mp}

\begin{leanbox}
inductive MPAction where $\mid$ heads $\mid$ tails\\[3pt]
def matchingPennies : NFGGame Bool (fun \_ $\Rightarrow$ MPAction) where\\
\bd Outcome := $\forall$ \_ : Bool, MPAction\\
\bd outcome := id\\
\bd utility := -- (H,H) = (1,-1), (H,T) = (-1,1), ...
\end{leanbox}

\paragraph{Faithfulness check.}\
There is no pure Nash equilibrium:\ at $(H, H)$ player $F$ deviates
to $T$ and gains $+2$; at $(H, T)$ player $T$ deviates to $T$ and
gains $+2$; the four pure profiles are similarly checked.  The
textbook mixed equilibrium is the symmetric uniform profile
$\Profile^*_T = \Profile^*_F = \tfrac12 \Dirac_H + \tfrac12 \Dirac_T$,
at which both players' EU is $0$.  The Lean example records
the no-pure-Nash result; it does not package the mixed-equilibrium
calculation as a theorem in \TT{Languages/NFG/Examples.lean}.

\subsection{A two-step centipede in EFG}
\label{sec:supp:examples:centipede}

A two-step centipede:\ player 0 chooses Take or Pass; if Pass,
player 1 chooses Take or Pass; payoffs $\mathrm{TT_0} = (1, 0)$,
$\mathrm{PT} = (0, 2)$, $\mathrm{PP} = (3, 1)$.  The textbook
backward-induction solution is $(T, T)$:\ at the second node,
player 1 strictly prefers $T$ ($2 > 1$), so they play $T$; at the
first node, player 0 anticipates this and prefers $T$ ($1 > 0$).
The library's packaged extensive-form example is the entry-deterrence
game in \TT{Languages/EFG/Refinements.lean}, where it proves:
\begin{itemize}\itemsep1pt
\item \TT{entrySPE_isSPE}:\ the backward-induction profile is SPE.
\item \TT{entryNash_isNash}:\ another profile is Nash.
\item \TT{entryNash_not_spe}:\ that Nash profile is not SPE.
\end{itemize}
The two-step centipede calculation is illustrative prose, not a
separate Lean theorem.

\subsection{A small Bayesian game}
\label{sec:supp:examples:bayesian}

A two-player coordination game with binary types:\ each player has
type $\theta_i \in \{0, 1\}$ drawn independently uniform; actions
are $a_i \in \{L, R\}$; payoffs are $u_i(\theta, a) = [\![ a_T = a_F ]\!]
\cdot [\![ a_T = \theta_T ]\!]$.  The Bayes-Nash equilibrium has
each player play $a_i = \theta_i$;\ the Lean proof unfolds
\TT{bayesNash_iff_exAnteEU} and verifies the four-deviation
inequality directly under the uniform prior.

% =====================================================================
\section{Reading order}
\label{sec:supp:reading}

Three suggested traversals.

\paragraph{(a)~"I want to use the library on a finite NFG."}\
\S\ref{sec:supp:probability} (probability primitives) →
\S\ref{sec:supp:gameform}, \S\ref{sec:supp:kernelgame},
\S\ref{sec:supp:eu} (semantic core) →
\S\ref{sec:supp:isNash} (Nash) →
\S\ref{sec:supp:nfg} (NFG) →
\S\ref{sec:supp:nash-exists} (existence).
\S\ref{sec:supp:examples:pd} is the canonical worked example.

\paragraph{(b)~"I want to add a new language."}\
\S\ref{sec:supp:core} (semantic core) →
\S\ref{sec:supp:nfg} or \S\ref{sec:supp:efg} as a template →
\S\ref{sec:supp:expressiveness} (Language/Reduction abstraction)
to register the language for bridges.  The
worked subclass files inside \TT{Languages/NFG/} (CongestionGame,
Bargaining, Stackelberg) are useful as additional templates.

\paragraph{(c)~"I want to follow the Kuhn theorem proof end-to-end."}\
\S\ref{sec:supp:probability} → \S\ref{sec:supp:gameform},
\S\ref{sec:supp:kernelgame} → \S\ref{sec:supp:efg} (EFG behavioral
profiles) → \S\ref{sec:supp:kuhn} (the four hypotheses, then the
two anchor theorems).  The proof split is in
\TT{BehavioralToMixedCore.lean} and \TT{MixedToBehavioralCore.lean}
in \TT{Theorems/Kuhn/}; the
\TT{CorrelatedRealization.lean} alternative goes via the
mixed-Nash-induced correlated distribution.

\bibliographystyle{plainnat}
\bibliography{references}

\end{document}
